\documentclass[UTF8,a4paper,11pt,openany]{ctexbook}
\usepackage{../common/statlearnbook}
\SLSetBookMetadata{第 5 册}{统计学习方法讲义 v2.7.0}{采样方法、主题模型与图排序}
\begin{document}
\setcounter{page}{576}
\setcounter{chapter}{31}
\setcounter{figure}{3}
\pagestyle{headings}
\chaptermark{蒙特卡罗方法与直接采样}
\sectionmark{接受--拒绝抽样}
图\ref{fig:V5-C02-rejection-envelope}取 $p(y)=6y(1-y)$、$q(y)=1$（$y\in[0,1]$）与 $c=1.6=8/5$。
此时 $p\le cq$ 几乎处处、最小包络差为 $1/10$；随机点 $(Y,Ucq(Y))$ 满足 $Ucq(Y)\le p(Y)$ 时接受。
接受率为 $5/8$，每个接受值所需平均提议数为 $8/5$；圆点 $(1/4,4/5)$ 接受，三角点 $(3/4,27/20)$ 普通拒绝但包络仍合法。
% v2.7.0 figure UID: FIG-P577-01; legacy UID: FIG-V5-C02-05
% FIGURE-ALT-COMBINED: 对象：在支撑集[0,1]上，实线为p(y)=6y(1-y)，虚线为cq(y)=8/5，浅填充为两者间的拒绝区；关系：p不超过cq（几乎处处），随机点(Y,Ucq(Y))落在p曲线下方或边界上时接受，否则普通拒绝；数值：最小包络差1/10，接受率5/8、平均提议数8/5、拒绝区面积3/5，圆点(1/4,4/5)接受而三角点(3/4,27/20)拒绝；结论：第二点仍低于合法包络，普通拒绝不等于包络失效。
\pgfkeysifdefined{/tikz/alt}{}{\tikzset{alt/.code={}}}
\tikzset{slfig-FIG-P577-01/.style={font=\fontsize{9.6pt}{11.6pt}\selectfont,every path/.append style={line cap=round,line join=round}}}
\pgfplotsset{slfig-FIG-P577-01-axis/.style={
  tick label style={font=\fontsize{9.6pt}{11.6pt}\selectfont},
  label style={font=\fontsize{9.6pt}{11.6pt}\selectfont},
  axis line style={line width=.72pt},tick style={line width=.60pt},clip=false}}
\begin{figure}[htbp]
\centering
\begin{tikzpicture}[slfig-FIG-P577-01,
  alt={对象：支撑集[0,1]上的实线p(y)=6y(1-y)、虚线cq(y)=8/5与浅填充拒绝区；关系：p不超过cq（几乎处处），最小差为1/10；圆点y=1/4、h=4/5、p=9/8、U=1/2不超过45/64，故接受；三角点y=3/4、h=27/20、p=9/8、U=27/32大于45/64，故普通拒绝但包络仍合法；接受率为5/8，平均提议数为8/5，拒绝区面积为3/5；结论：普通拒绝不等于包络失效。}]
\begin{axis}[slfig-FIG-P577-01-axis,width=.88\textwidth,height=8.8cm,
  xmin=-.035,xmax=1.035,ymin=0,ymax=1.78,
  axis lines=left,xlabel={$y$},ylabel={密度高度},
  xtick={0,.25,.5,.75,1},xticklabels={$0$,$\frac14$,$\frac12$,$\frac34$,$1$},
  ytick={0,.4,.8,1.2,1.6},
  xlabel style={yshift=-7pt},ylabel style={yshift=3pt},clip=false]
  \addplot[name path=p,draw=SLBlue,line width=1.15pt,domain=0:1,samples=301]{6*x*(1-x)};
  \addplot[name path=cq,draw=SLTeal,line width=.90pt,densely dashed,domain=0:1,samples=301]{1.6};
  \addplot[fill=SLTeal,fill opacity=.09,draw=none] fill between[of=cq and p];

  \node[anchor=south,draw=SLGray,rounded corners=1.5pt,fill=white,
    line width=.55pt,inner xsep=3mm,inner ysep=1.4mm,text width=125mm]
    at (axis description cs:.5,1.055)
    {{\bfseries\fontsize{10.2pt}{12.2pt}\selectfont 合法包络与含边界接受门}\\[-.2mm]
     $p(y)=6y(1-y),\ q(y)=1,\ c=\frac85=1.6\quad(y\in[0,1])$\\
     $p\le cq\ \text{（几乎处处）},\quad
      \min_{[0,1]}(cq-p)=\frac1{10}\ \text{（在 }y=\frac12\text{）}$\\
     $Ucq(Y)\le p(Y)\Longleftrightarrow\text{接受（含边界）};\quad
      P(\text{接受})=\frac58,\ \text{平均提议数/接受值}=\frac85,\
      \int_0^1(cq-p)\,\mathrm dy=\frac35$};

  \node[anchor=south east,text=SLBlue,fill=white,inner sep=.8pt]
    at (axis cs:.40,1.39) {实线 $p(y)$};
  \node[anchor=south east,text=SLTeal,fill=white,inner sep=.8pt]
    at (axis cs:.97,1.60) {虚线 $cq(y)=\frac85$};
  \draw[<->,draw=SLTextGray,line width=.72pt]
    (axis cs:.50,1.50)--(axis cs:.50,1.60)
    node[midway,right=1.5pt,fill=white,inner sep=.5pt] {$\frac1{10}$};
  \node[anchor=west,text=SLTextGray,fill=white,inner sep=1pt,align=left]
    at (axis cs:.055,1.15) {浅填充：包络差\\拒绝区面积 $\frac35$};

  \draw[draw=SLTeal,line width=.74pt] (axis cs:.25,0)--(axis cs:.25,.80);
  \addplot[only marks,mark=*,mark size=3.1pt,draw=SLTeal,fill=SLTeal]
    coordinates {(.25,.80)};
  \node[anchor=north east,draw=SLTeal,rounded corners=1.2pt,fill=white,
    line width=.55pt,inner sep=1.2mm,align=right]
    at (axis cs:.235,.77)
    {接受（圆点）\\$y=\frac14,\ h=\frac45,\ p=\frac98$\\
     $U=\frac{h}{cq}=\frac12\le\frac{p}{cq}=\frac{45}{64}$};

  \draw[draw=SLGold,line width=.74pt,densely dotted] (axis cs:.75,0)--(axis cs:.75,1.35);
  \addplot[only marks,mark=triangle*,mark size=3.6pt,draw=SLGold,fill=white,
    line width=.85pt] coordinates {(.75,1.35)};
  \node[anchor=north west,draw=SLGold,rounded corners=1.2pt,fill=white,
    line width=.55pt,inner sep=1.2mm,align=left]
    at (axis cs:.750,1.32)
    {普通拒绝（三角点；包络合法）\\
     $y=\frac34,\ h=\frac{27}{20},\ p=\frac98$\\
     $U=\frac{h}{cq}=\frac{27}{32}>\frac{p}{cq}=\frac{45}{64}$};

  \addplot[only marks,mark=square*,mark size=2.4pt,draw=SLTextGray,fill=white]
    coordinates {(0,0) (1,0)};
  \node[anchor=north west,text=SLTextGray,yshift=-3pt]
    at (axis cs:0,0) {支撑集边界};
  \node[anchor=north east,text=SLTextGray,yshift=-3pt]
    at (axis cs:1,0) {支撑集边界};
\end{axis}
\end{tikzpicture}
\caption{包络满足 $p\le cq$（几乎处处）；随机点 $(Y,Ucq(Y))$ 落在 $p$ 曲线下方（含边界）时接受 $Y$。}\label{fig:V5-C02-rejection-envelope}
\end{figure}

\FloatBarrier
\noindent\textbf{读图检查。}先看实线 $p$ 与虚线 $cq$ 的合法支配和最小差；
再按含边界门 $Ucq(Y)\le p(Y)$ 核对圆点与三角点；
最终由面积确认接受率 $5/8$、平均提议数 $8/5$ 与拒绝区面积 $3/5$，并区分普通拒绝和包络失效。
\begin{comparisonbox}
\textbf{包络条件是算法合法性的前提。} 若某处 $p(y)>cq(y)$，接受比率会越过 1；此时必须停止并修正包络，不能把它当作一次普通拒绝。
\end{comparisonbox}
\end{document}
